\section{Integral and finite fibres over completed residue fields}
\label{sec:general-finite-fibres}

\begin{theorem}[Extension of a complete ultrametric norm]
\label{spectrum:complete-norm-extension}
\lean{EnvelopeProperty.Valuation.le_one_of_trivial_of_isAlgebraic,
EnvelopeProperty.Valuation.eq_one_of_trivial_of_isAlgebraic,
EnvelopeProperty.Valuation.exists_extends_norm,
EnvelopeProperty.Valuation.eq_of_extends_norm}
\leanok
Let $K$ be a complete ultrametric normed field and $L/K$ an algebraic
field extension. There exists exactly one real-valued ultrametric
valuation on $L$ restricting to the norm of $K$, including when
that norm is trivial.
\end{theorem}
\begin{proof}
For a nontrivial norm on $K$, the spectral norm on the algebraic
extension is multiplicative and ultrametric and extends the base
norm. The uniqueness theorem for extensions over a complete field
identifies every extending field absolute value with this spectral
norm. For a trivial base norm, apply a monic equation to bound the
valuation of every algebraic element by one. Apply the same bound
to the inverse of every nonzero element. Multiplicativity then
forces every nonzero value to equal one. The trivial valuation
therefore supplies the unique extension in this case.
\end{proof}

\begin{theorem}[Integral spectra over a complete normed field]
\label{spectrum:integral-field-support-classification}
\lean{EnvelopeProperty.Berkovich.RelativeSpectrum.isMaximal_supp_of_isIntegral,
EnvelopeProperty.Berkovich.RelativeSpectrum.eq_of_supp_eq_of_isIntegral,
EnvelopeProperty.Berkovich.RelativeSpectrum.integralSupport,
EnvelopeProperty.Berkovich.RelativeSpectrum.integralSupport_surjective,
EnvelopeProperty.Berkovich.RelativeSpectrum.integralSupport_bijective,
EnvelopeProperty.Berkovich.RelativeSpectrum.integralMaximalSpectrumEquiv}
\leanok
\uses{spectrum:complete-norm-extension,spectrum:quotient-homeomorphism}
If a commutative $K$-algebra $C$ is integral over a complete
ultrametric normed field $K$, its spectrum with prescribed base norm
is in bijection with $\operatorname{MaxSpec}(C)$ by its actual
valuation support. No finiteness or Noetherian assumption is needed.
\end{theorem}
\begin{proof}
Each valuation kernel is prime. Its contraction to $K$ is the zero
prime, which is maximal; integrality makes the kernel maximal in
$C$. For a maximal ideal $\mathfrak m$, the quotient field
$C/\mathfrak m$ is algebraic over $K$. Extend the complete base norm
to this actual field and pull back along the quotient map. This
gives a prescribed-base seminorm with support precisely $\mathfrak m$.
If two points have the same support, they descend to valuations on
the same quotient field, both extending the base norm. Uniqueness
of the norm extension identifies them on the whole quotient, hence
on every original ring element. This proves bijectivity.
\end{proof}

\begin{theorem}[Completed tensor classification of integral fibres]
\label{spectrum:integral-completed-fibres}
\lean{EnvelopeProperty.Berkovich.Spectrum.fibrePoint_supp,
EnvelopeProperty.Berkovich.Spectrum.isMaximal_ker_fibreHom,
EnvelopeProperty.Berkovich.Spectrum.fibreToMaximalSpectrum,
EnvelopeProperty.Berkovich.Spectrum.fibreToMaximalSpectrum_bijective,
EnvelopeProperty.Berkovich.Spectrum.fibreEquivMaximalSpectrum,
EnvelopeProperty.Berkovich.Spectrum.eq_of_ker_fibreHom_eq}
\leanok
\uses{spectrum:integral-field-support-classification,spectrum:fibre-equivalence}
Let $A\to B$ be an integral homomorphism of arbitrary commutative
rings and $x\in\mathcal S(A)$. The actual fibre over $x$ is in
bijection with $\operatorname{MaxSpec}(\mathcal H(x)\otimes_A B)$.
A fibre point maps to the kernel of its genuine tensor
homomorphism $\mathcal H(x)\otimes_A B\to\mathcal H(y)$.
\end{theorem}
\begin{proof}
The genuine tensor fibre equivalence identifies the fibre with the
spectrum of $\mathcal H(x)\otimes_A B$ having the prescribed norm
of $\mathcal H(x)$. This tensor algebra is integral over
$\mathcal H(x)$ by integral base change. Its coefficient field is
complete and ultrametric, so the preceding classification applies.
The tensor valuation is the pullback of the norm on $\mathcal H(y)$;
its support is therefore exactly the kernel of the actual tensor
homomorphism. This identifies the resulting bijection as asserted.
\end{proof}

\begin{corollary}[Finite fibres]
\label{spectrum:finite-completed-fibres}
\lean{EnvelopeProperty.Berkovich.RelativeSpectrum.finite_normSpectrum,
EnvelopeProperty.Berkovich.Spectrum.finite_fibre,
EnvelopeProperty.Berkovich.Spectrum.finite_preimage_comap_singleton}
\leanok
\uses{spectrum:integral-completed-fibres}
If $B$ is finite as an $A$-module, each seminorm fibre is a finite set.
\end{corollary}
\begin{proof}
The tensor fibre algebra is finite dimensional over $\mathcal H(x)$,
hence Artinian. The maximal spectrum of an Artinian ring is finite.
The actual fibre bijection therefore gives finiteness of the fibre.
The argument includes the zero tensor algebra and empty fibres.
\end{proof}

\begin{corollary}[Arbitrary prescribed-base fibres]
\label{spectrum:relative-integral-completed-fibres}
\lean{EnvelopeProperty.Berkovich.RelativeSpectrum.Fibre,
EnvelopeProperty.Berkovich.RelativeSpectrum.fibreForgetEquiv,
EnvelopeProperty.Berkovich.RelativeSpectrum.fibreEquivMaximalSpectrum,
EnvelopeProperty.Berkovich.RelativeSpectrum.finite_fibre,
EnvelopeProperty.Berkovich.RelativeSpectrum.finite_preimage_comap_singleton}
\leanok
\uses{spectrum:integral-completed-fibres,spectrum:finite-completed-fibres,spectrum:relative}
For compatible scalar towers $R\to A\to B$ and any prescribed
seminorm $v$ on $R$, an integral morphism's fibre over
$x\in\mathcal S_v(A)$ has the same completed-tensor classification.
For finite algebra homomorphisms these fibres are finite.
\end{corollary}
\begin{proof}
Forgetting the prescribed base gives a point of the corresponding
absolute fibre. Conversely an absolute point over $x$ automatically
has the required base values: evaluate the scalar tower on every
$r\in R$, use its equality over $x$, then use the defining base
restriction of $x$. These are inverse maps on the entire seminorm.
Compose with the absolute fibre classification. An arbitrary ring
or algebra homomorphism is handled by its actual induced algebra
structure, preserving the stated integral or finite hypothesis.
\end{proof}

No classification hypothesis is assumed. Integrality is used for
the maximal-support bijection, and module finiteness is used only
for fibre finiteness. These statements do not assert finite
cardinality for general integral fibres.
